\section{Síntesis y ampliación}
\begin{table}[H]
\centering
\begin{tabular}{p{4.5cm}p{11cm}}
\toprule
Tema & Puntos clave de la clase \\
\midrule
Arquitectura de autorrecompensa & El Task Generator genera tareas diversas, el Solution Generator produce Chain-of-Thought, el Reward Model se autoevalúa mediante verified response. \\
Self-Instruct + Verified Response & Self-Instruct asegura la amplitud de las tareas, verified response aporta muestras de cross-check y evita el encadenamiento excesivo. \\
Meta-Judge y evaluación & Meta-Judge usa SelfT evaluators para simular al juez, logra distinguir entre draft y verified response y supera de forma efectiva el plateau. \\
Retos de despliegue & el inference-time compute debe ser controlable, la alucinación requiere multi-evaluator y un sistema self-aware puede solicitar la intervención humana de forma más razonable. \\
\bottomrule
\end{tabular}
\caption{Hilo conductor central del self-improving pipeline visto en esta clase.}
\end{table}

\begin{table}[H]
\centering
\begin{tabular}{p{5cm}p{10cm}}
\toprule
Dirección de investigación & Descripción \\
\midrule
Recompensa verificable & Construir verified response para que el modelo de recompensa mantenga su confiabilidad sin necesidad de anotación humana. \\
SelfT + Meta Judge & Usar SelfT evaluators para entrenar al meta judge, recuperando de forma automática un juicio de alta calidad a partir del draft y el verified response. \\
Interacción en línea & Mediante el mecanismo Self-Feeding se enlazan la interacción real, el verified response y los datos de instrucción, formando un ciclo de aprendizaje continuo. \\
\bottomrule
\end{tabular}
\caption{Hoja de ruta de investigación a la que conviene dar seguimiento continuo.}
\end{table}

\subsection{Lecturas adicionales}
\begin{itemize}
    \item El artículo de OpenAI «Self-Feeding Chatbot», que muestra la iteración de verified response en el diálogo.
    \item La investigación de SelfT Evaluators / Self-Rewarding de Weston y colegas, que explica en detalle el entrenamiento del juez.
    \item El benchmark Alpaca Eval 2, que ofrece una referencia de win rate junto con la verificación mediante meta judge.
    \item Artículos comparativos entre Chain-of-Thought y RLHF/DPO, útiles para decidir cuándo usar self-rewarding.
\end{itemize}

\subsection{Resumen de la sección}
Esta clase de Jason Weston nos muestra que la capacidad de razonamiento es una cadena que necesita tres pasos: plantear la pregunta, verificar la respuesta y evaluar el resultado. Solo cuando el modelo aprende por sí mismo a partir del verified response y vuelve a confirmarlo bajo la guía del meta judge, el entrenamiento de reasoning puede avanzar al siguiente nivel.

\end{document}
